% Document class and basic packages
\documentclass[a4paper]{article}

% Essential packages
\usepackage{fullpage}
\usepackage{amsmath}
\usepackage{amssymb}
\usepackage{textcomp}
\usepackage[utf8]{inputenc}
\usepackage[T1]{fontenc}
\usepackage[hidelinks]{hyperref}
\usepackage[left=2cm, right=2cm, top=2cm]{geometry}
\usepackage{longtable}

% Document settings
\textheight=10in
\pagestyle{empty}
\raggedright

% Custom commands
\def\bull{\vrule height 0.8ex width .7ex depth -.1ex}

\newcommand{\area}[2]{
    \vspace*{-9pt}
    \begin{verse}
        \textbf{#1}   #2
    \end{verse}
}

\newcommand{\lineunder}{
    \vspace*{-8pt} \\
    \hspace*{-18pt} \hrulefill \\
}

\newcommand{\header}[1]{
    {\hspace*{-18pt}\vspace*{6pt} \textsc{#1}}
    \vspace*{-6pt} \lineunder
}

\newcommand{\employer}[3]{
    { \textbf{#1} (#2)\\ \underline{\textbf{\emph{#3}}}\\  }
}

\newcommand{\contact}[3]{
    \vspace*{-10pt}
    \begin{center}
        {\Huge \scshape {#1}}\\
        #2 \\ #3
    \end{center}
    \vspace*{-8pt}
}

\newenvironment{achievements}{
    \begin{list}
        {$\bullet$}{\topsep 0pt \itemsep -2pt}}{\vspace*{4pt}
    \end{list}
}

\newcommand{\schoolwithcourses}[4]{
    \textbf{#1} #2 $\bullet$ #3\\
    #4 \\
    \vspace*{5pt}
}

\newcommand{\school}[4]{
    \textbf{#1} #2 $\bullet$ #3\\
    #4 \\
}


% Begin document
\begin{document}
    \vspace*{-40pt}

    % Header section
    \vspace*{-2pt}
    \begin{center}
        {\Huge \scshape {Lucas Marujo}}\\
        \vspace*{2pt}
        \ {Desenvolvedor Full-Stack}\\
        \vspace*{2pt}
        \href{mailto:lucas.m.amadeu@gmail.com}{lucas.m.amadeu@gmail.com} | \href{tel:+55 11 992515086}{+55 11 992515086}\\
        \vspace*{4pt}
        \textbf{\href{https://www.linkedin.com/in/lucas-amadeu/}{Linkedin | }}\textbf{\href{http://github.com/lucasmarujo}{GitHub | }}\textbf{\href{https://marujo.dev/}{Portfolio}}\\
    \end{center}

    % Seção de Resumo
    \header{Resumo}
    \vspace{2mm}
    Engenheiro de Software Full Stack especializado em Python, Django, React, React Native e aplicações com Inteligência Artificial. Experiência no desenvolvimento de sistemas backend escaláveis, aplicações web e mobile modernas, e soluções baseadas em LLM utilizando LangChain. Apaixonado por construir software pronto para produção, infraestrutura em nuvem e automação inteligente.
    \vspace{2mm}

    % Seção de Experiência
    \header{Experiência}
    \vspace{2mm}

\textbf{GreenLegis} | \textbf{Engenheiro de Software} \hfill Fev 2026 – Atual
\vspace{-3mm}
\begin{itemize} \itemsep -3pt
    \item Desenvolvi e mantive \textbf{soluções baseadas em IA} utilizando \textbf{Python}, \textbf{LangChain} e ecossistemas modernos de LLMs, criando agentes inteligentes e fluxos de trabalho prontos para produção.
    \item Projetei e implementei \textbf{agentes de IA}, estratégias de engenharia de prompts, pipelines de recuperação de contexto (RAG) e integrações com ferramentas para automatizar processos jurídicos e de negócio.
    \item Integrei \textbf{serviços de IA} em aplicações legadas desenvolvidas em \textbf{.NET}, adicionando funcionalidades inteligentes sem comprometer a arquitetura e a estabilidade dos sistemas existentes.
    \item Trabalhei com \textbf{PostgreSQL}, \textbf{MySQL} e \textbf{MongoDB}, projetando e mantendo modelos de dados para aplicações de IA e sistemas corporativos.
    \item Colaborei com equipes multidisciplinares para entregar funcionalidades escaláveis de IA, APIs REST e integrações de serviços em um ambiente de desenvolvimento ágil.
    \item Contribuí para a arquitetura e evolução de produtos orientados por IA, com foco em escalabilidade, manutenibilidade e confiabilidade em produção.
\end{itemize}

\textbf{Grupo Tombini} | \textbf{Engenheiro de Software} \hfill Abr 2025 – Fev 2026
\vspace{-3mm}
\begin{itemize} \itemsep -3pt
    \item Liderei a migração da infraestrutura da empresa de um ambiente \textbf{on-premises} para a \textbf{AWS}, aumentando a escalabilidade, confiabilidade e eficiência dos processos de implantação.
    \item Projetei e desenvolvi uma \textbf{aplicação de onboarding para motoristas} utilizando \textbf{Django} e \textbf{React Native} (iOS e Android), atualmente utilizada por mais de \textbf{2.500 usuários ativos}.
    \item Otimizei o ERP principal da empresa, reduzindo o tempo de resposta de aproximadamente \textbf{40 segundos} para \textbf{1,6 segundos} por meio de otimizações no backend, banco de dados e consultas.
    \item Desenvolvi aplicações web e mobile de ponta a ponta utilizando \textbf{Python (Django)}, \textbf{TypeScript}, \textbf{React.js}, \textbf{Node.js} e \textbf{React Native}.
    \item Implementei \textbf{soluções baseadas em IA} utilizando \textbf{LangChain}, modelos YOLO e automações de processos com \textbf{n8n} e \textbf{Selenium}, eliminando tarefas manuais e aumentando a eficiência operacional.
    \item Liderei migrações de bancos de dados em larga escala, modernizei módulos legados e contribuí para a arquitetura dos sistemas, engenharia de processos e implantações utilizando \textbf{Docker}.
\end{itemize}

    \textbf{Rede Boa Supermercados}\textbf{ | Estagiário em Suporte de TI } \hfill  Out 2022 - Jan 2023\\
    \vspace{-3mm}
    \begin{itemize} \itemsep -3pt
        \item Manutenção de redes, hardware e sistemas de PDV, Suporte técnico n1 e n2 em ERP Protheus e Linux, resolução de falhas no sistema.
    \end{itemize}

    \textbf{AirBlower Engenharia}\textbf{ | Web Developer } \hfill  Fev 2020 - Dez 2021\\
    \vspace{-3mm}
    \begin{itemize} \itemsep -3pt
        \item Participação no design, desenvolvimento e deploy do website institucional da empresa utilizando \textbf{HTML}, \textbf{CSS} e \textbf{JavaScript}, além da criação de um sistema web para cálculo e geração de orçamentos de ventiladores industriais. Também atuei na administração das redes sociais e na manutenção de infraestrutura de redes, hardware e software.
    \end{itemize}

    \textbf{Fly Wheel (Intellibrand)}\textbf{ | Estagiário em TI } \hfill  Jun 2019 - Jan 2020\\
    \vspace{-3mm}
    \begin{itemize} \itemsep -3pt
        \item Monitoramento de plataforma digital interna.
    \end{itemize}

        % Projects section
    \header{Projetos}
    \vspace{2mm}
    
    \textbf{Pulso}\text{ | Plataforma de treino, dieta e IA para Android, iOS e Web} \hfill  2025 -- Atual\\
    \vspace{-3mm}
    \begin{itemize} \itemsep -3pt
        \item Desenvolvimento de uma plataforma completa utilizando \textbf{React Native}, \textbf{React.js}, \textbf{Python}, \textbf{Django} e \textbf{LangChain}, composta por um aplicativo para \textbf{Android} e \textbf{iOS} voltado ao controle de treinos, dieta e acompanhamento físico com IA integrada.
        \item Desenvolvimento de um \textbf{dashboard para personal trainers}, permitindo o gerenciamento de alunos, acompanhamento de evolução, planos de treino, métricas e funcionalidades baseadas em Inteligência Artificial.
    \end{itemize}
    
    \textbf{Orkai}\text{ | Infinite Canvas para orquestração de agentes de IA} \hfill  2025 -- Atual\\
    \vspace{-3mm}
    \begin{itemize} \itemsep -3pt
        \item Desenvolvimento de uma plataforma de \textbf{Infinite Canvas} para criação e orquestração de fluxos de agentes de IA, permitindo a construção visual de pipelines inteligentes e automações complexas.
        \item Implementação de arquitetura baseada em \textbf{Python}, \textbf{LangChain}, \textbf{OpenRouter}, \textbf{Docker} e modelos de linguagem modernos para coordenação, comunicação e execução de agentes autônomos.
    \end{itemize}
    

    % Skills section
    \header{Habilidades}
    \vspace{2mm}
    \begin{longtable}{p{4cm}p{12cm}}
        Frontend: & React.js, React Native, Next.js, Typescript, TailwindCSS, HTML, CSS, Javascript. \\
        Backend: & C\#, .NET, Python, Django, Golang, Node.js, Nest. \\
        Banco de dados \& Servers: & PostgreSQL, MySQL, MongoDB. \\
        Plataformas \& Outros: &  AWS, Azure, Git, Github, Gitlab, BitBucket, Jira, Jenkins, Linux, CICD. \\
    \end{longtable}
    \vspace{1mm}

    % Education section
    \header{Educação}
    \vspace{2mm}
    \textbf{FIAP}\hfill São Paulo, São Paulo, Brasil\\
    MBA EM AI ENGINEER E MULTI-AGENTS (Cursando) \hfill Abr 2026 - Fev 2027\\
    \vspace{2mm}
    \textbf{Universidade Padre Anchieta}\hfill Jundiaí, São Paulo, Brasil\\
    BACHARELADO EM CIÊNCIAS DA COMPUTAÇÃO (Completo) \hfill Jan 2022 - Dez 2025\\
    \vspace{2mm}
    \textbf{ILAC - International Language Academy of Canada}\hfill Toronto, Canada\\
    ENGLISH AS SECOND LANGUAGE (Completo) - CPE C2 level \hfill Nov 2024 - Dez 2024\\
    \vspace{2mm}
    \textbf{Faculdade Flamingo}\hfill Lapa, São Paulo, Brasil\\
    TÉCNICO EM INFORMÁTICA(Completo) \hfill Jan 2019 - Dez 2021\\

\end{document}